\subsubsection{Calcul par fronts sur les révisions d’un haplotype}
Les révisions v1.1 et v1.2 de HG002 permettent de mesurer l’effet de quelques éditions dispersées sur un chromosome entier~\cite{hg002,hg002data}. Chaque comparaison porte sur deux versions du même haplotype. Les séquences du chromosome 1 comptent de 244 à 252 millions de bases. Sur le chromosome 21, l'assemblage maternel conserve 3\,671\,963 symboles N et le paternel 2\,300\,690. Ces lacunes d’assemblage restent dans le calcul, avec un coût nul pour N face à N. L’application des 67 entrées du VCF aux quatre séquences v1.1 restitue les quatre séquences v1.2.

Dans les fragments de la figure~\ref{fig:split}, les longues plages identiques peuvent se parcourir directement. Pour un budget d’éditions \(s\) et une diagonale \(k=i-j\), le front \(F_s(k)\) mémorise le plus grand indice \(i\) atteint. La délétion, la substitution et l’insertion donnent les trois candidats~\cite{ukkonen,wfa},
\[
 x=\max\{F_{s-1}(k-1)+1,\;F_{s-1}(k)+1,\;F_{s-1}(k+1)\},
 \qquad F_s(k)=x+\operatorname{LCP}(A_{x:},B_{x-k:}).
\]
Les transitions sortant du rectangle sont écartées. Le préfixe commun initialise \(F_0(0)\). Le premier front qui atteint \((m,n)\) donne la distance minimale. Pour le score seul, la commande \texttt{nvpd-sparse} conserve les coupes vérifiées et NvPD pour les petits rectangles. Dans les grands, elle explore jusqu’à 128 éditions par fragment avant de passer à BiWFA. Deux tableaux de 259 entiers occupent 2\,072 octets par travailleur. L’extension lit les deux séquences par blocs de 4\,096 puis de 64 octets.

Le panneau a) de la figure~\ref{fig:hg002-dispatch} compare les solveurs sur un cœur EPYC. Dans chacun des cinq blocs, un processus par méthode effectue un appel de chauffe puis un appel chronométré, après lecture des séquences.

Lorsqu’une coupe est rejetée, les fragments adjacents doivent être réunis. Les bases déjà reconnues comme identiques aux deux extrémités du groupe restent appariables. Pour des chaînes \(P,X,Y,S\), la distance unitaire vérifie
\[
 d(PXS,PYS)=d(X,Y).
\]
Le programme conserve la longueur des extrémités communes de chaque fragment. Il enlève les fragments terminaux de score nul, puis les extrémités communes des fragments restants. S’il ne reste qu’un fragment, son score est déjà connu. Le panneau b) isole l’effet de cette réutilisation à nombre de cœurs fixé. Pour chaque chromosome, un même processus exécute un bloc de chauffe puis sept blocs en ordre aléatoire.
\chart{comparison/chromosome_dispatch.pdf}{Le panneau a) rapporte le temps de WFA2 à celui des solveurs sur un cœur, avec l’écart-type et l’IC à 95\,\% de cinq rapports appariés. Le panneau b) rapporte le temps avec relecture à celui avec réutilisation, pour un et huit cœurs. Le symbole indique la moyenne, le trait épais un écart-type et le trait fin les extrêmes des sept blocs.}{fig:hg002-dispatch}{python/studies/chromosome_dispatch.py}
Sur le chromosome 1 maternel, le score vaut 54. À un cœur, le temps passe de \(157{,}05\pm0{,}47\) à \(28{,}19\pm0{,}06\)\,ms lorsque le découpage utilise le front borné. WFA2 prend \(87{,}18\pm0{,}80\)\,ms et Edlib \(4{,}393\pm0{,}031\)\,s. Le pic résident moyen du processus descend de 834 à 477\,Mio entre les deux variantes du découpage, FASTA compris. Les compteurs en mode utilisateur relèvent 59,6 millions d’instructions contre 334,3 millions. Le nombre de défauts de lecture L1 passe de 38,89 à 11,03 millions, alors que leur part dans les lectures monte de 23,3 à 35,6\,\%.

Sur le chromosome~1 paternel, une des sept coupes proposées entraîne une fusion. À huit cœurs, la réutilisation ramène le temps de \(21{,}07\) à \(10{,}76\)\,ms. Sur le chromosome~1 maternel, une coupe est également rejetée et le gain moyen vaut 1,054. Sur le chromosome~21, toutes les coupes sont acceptées pour les deux haplotypes. Les rapports moyens à huit cœurs valent alors 0,912 et 0,956 pour les haplotypes maternel et paternel. Leur étendue comprend 1 dans les deux cas.

\synthese{À huit cœurs EPYC, le gain lié à la réutilisation des extrémités atteint \(1{,}972\pm0{,}182\) sur le chromosome~1 paternel de HG002. Les sept rapports mesurés dans le même processus vont de 1,637 à 2,201 pour un score de 111 éditions.}
